\chapter{Complexity — Finite}
%%% generated by the blueprint pipeline from TCSlib.Complexity.TuringMachine.Finite
%%%%%%%%%%%%%%%%%%%%%%%%%%%%%%%%%%%%%%%%%%%%%%%%%%%%%%%%%%%%%%%%%%%%%%%%%%%%%%%
\section{Overview}

This module bundles a deterministic multi-tape Turing machine with a finite
state set, giving the machine model of Arora--Barak \S 1.2 over which all the
headline complexity-theoretic definitions are stated.  On top of the bundled
machines it defines what it means to compute an output within a time bound, to
compute a string function in time $T(n)$ (also across an alphabet embedding),
and to compute a function with no time constraint, and it proves the basic
laws of these notions: time bounds can be weakened, no machine halts in zero
steps, completed outputs are unique, and a total machine's completed outputs
are exactly its prescribed values.  Two raw-layer lemmas record that the
output tape gains at most one symbol per step and only ever grows.

%%%%%%%%%%%%%%%%%%%%%%%%%%%%%%%%%%%%%%%%%%%%%%%%%%%%%%%%%%%%%%%%%%%%%%%%%%%%%%%
\section{Declarations}

\begin{theorem}[Output length is at most the number of steps]
\lean{Turing.MultiTapeTM.output_length_le}
\label{Turing.MultiTapeTM.output_length_le}
\leanok
\uses{Turing.MultiTapeTM, Turing.MultiTapeTM.initCfg, Turing.MultiTapeTM.outputSymbol, Turing.MultiTapeTM.runFrom, Turing.MultiTapeTM.runFrom_succ_eq_step', Turing.MultiTapeTM.step_output}
\difficulty{4}
Let $M$ be a deterministic multi-tape Turing machine, let $x$ be an input
word, and let $t \in \mathbb{N}$.  Then the word written on the output tape
after $M$ runs for $t$ steps from its initial configuration on $x$ has length
at most $t$.
\end{theorem}

\begin{theorem}[Output is monotone along a run]
\lean{Turing.MultiTapeTM.output_prefix}
\label{Turing.MultiTapeTM.output_prefix}
\leanok
\uses{Turing.Cfg, Turing.MultiTapeTM, Turing.MultiTapeTM.runFrom, Turing.MultiTapeTM.runFrom_add, Turing.MultiTapeTM.runFrom_succ_eq_step', Turing.MultiTapeTM.step_output}
\difficulty{4}
Let $M$ be a deterministic multi-tape Turing machine, let $c$ be any
configuration of $M$, and let $t \le t'$ be natural numbers.  Then the word on
the output tape after running $M$ for $t$ steps from $c$ is a prefix of the
word on the output tape after running $M$ for $t'$ steps from $c$.
\end{theorem}

\begin{definition}[Bundled finite Turing machine]
\lean{Turing.FinTM}
\label{Turing.FinTM}
\leanok
\uses{Turing.MultiTapeTM}
For an alphabet $\Sigma$, a finite Turing machine consists of a number
$k \in \mathbb{N}$ of work tapes, a finite state set with decidable equality,
and a deterministic multi-tape Turing machine with $k$ work tapes and those
states, whose tape cells hold either a symbol of $\Sigma$ or a blank.  This is
the machine model of Arora--Barak \S 1.2, up to the declared model variations
(append-only output tape, start-marker-free initialization).
\end{definition}

\begin{definition}[Computing an output within a time bound]
\lean{Turing.FinTM.ComputesInTime}
\label{Turing.FinTM.ComputesInTime}
\leanok
\uses{Turing.FinTM, Turing.MultiTapeTM.ComputesInTimeAndSpace}
A finite Turing machine $M$ over an alphabet $\Sigma$ \emph{computes the
output $w$ from the input $x$ within $t$ steps} if there exists a space bound
$s$ such that $M$, started on $x$, halts within $t$ steps using space $s$ with
exactly $w$ written on its output tape; that is, the space used is
existentially quantified away from the time-and-space notion of computation
(Arora--Barak, Definition 1.3).
\end{definition}

\begin{definition}[Computing a function in time $T$]
\lean{Turing.FinTM.ComputesFunInTime}
\label{Turing.FinTM.ComputesFunInTime}
\leanok
\uses{Turing.FinTM, Turing.FinTM.ComputesInTime}
A finite Turing machine $M$ over an alphabet $\Sigma$ \emph{computes the
string function $f \colon \Sigma^{*} \to \Sigma^{*}$ in time
$T \colon \mathbb{N} \to \mathbb{N}$} if for every input word $x$, the machine
$M$ halts on $x$ within $T(\abs{x})$ steps with $f(x)$ written on its output
tape (Arora--Barak, Definition 1.3).
\end{definition}

\begin{definition}[Computing a function via an alphabet embedding]
\lean{Turing.FinTM.ComputesFunInTimeVia}
\label{Turing.FinTM.ComputesFunInTimeVia}
\leanok
\uses{Turing.FinTM, Turing.FinTM.ComputesInTime}
Let $\alpha$ and $\Gamma$ be alphabets, let $e \colon \alpha \hookrightarrow
\Gamma$ be an injection, and let $M$ be a finite Turing machine over $\Gamma$.
The machine $M$ \emph{computes the string function $f \colon \alpha^{*} \to
\alpha^{*}$ via $e$ in time $T \colon \mathbb{N} \to \mathbb{N}$} if for every
word $x$ over $\alpha$, on the input obtained by applying $e$ letterwise to
$x$, the machine $M$ halts within $T(\abs{x})$ steps with the letterwise image
of $f(x)$ under $e$ written on its output tape.  This is how a machine over a
larger alphabet is said to compute a function on a smaller one
(Arora--Barak \S 1.3.1).
\end{definition}

\begin{definition}[Computing a function, no time constraint]
\lean{Turing.FinTM.Computes}
\label{Turing.FinTM.Computes}
\leanok
\uses{Turing.FinTM, Turing.FinTM.ComputesInTime}
A finite Turing machine $M$ over an alphabet $\Sigma$ \emph{computes the
string function $f \colon \Sigma^{*} \to \Sigma^{*}$} if for every input word
$x$ there exists a number of steps $t$ such that $M$ halts on $x$ within $t$
steps with $f(x)$ written on its output tape.  This is the notion of
computability underlying the uncomputability results
(Arora--Barak \S 1.4--1.5).
\end{definition}

\begin{theorem}[Weakening the time bound]
\lean{Turing.FinTM.ComputesInTime.mono}
\label{Turing.FinTM.ComputesInTime.mono}
\leanok
\uses{Turing.FinTM, Turing.FinTM.ComputesInTime, Turing.MultiTapeTM.ComputesInTimeAndSpace, Turing.MultiTapeTM.initCfg, Turing.MultiTapeTM.runFrom, Turing.MultiTapeTM.runFrom_add, Turing.MultiTapeTM.runFrom_of_halt}
\difficulty{4}
Let $M$ be a finite Turing machine over an alphabet $\Sigma$, let $x$ and $w$
be words over $\Sigma$, and let $t \le t'$ be natural numbers.  If $M$,
started on $x$, halts within $t$ steps with $w$ on its output tape, then it
also halts within $t'$ steps with $w$ on its output tape.
\end{theorem}

\begin{theorem}[No computation in zero steps]
\lean{Turing.FinTM.not_computesInTime_zero}
\label{Turing.FinTM.not_computesInTime_zero}
\leanok
\uses{Turing.FinTM, Turing.FinTM.ComputesInTime, Turing.MultiTapeTM.runFrom_zero}
\difficulty{1}
For every finite Turing machine $M$ over an alphabet $\Sigma$ and all words
$x$ and $w$ over $\Sigma$, the machine $M$ does not compute the output $w$
from the input $x$ within $0$ steps: after zero steps the machine is still in
its initial, non-halted state.
\end{theorem}

\begin{theorem}[Uniqueness of the completed output]
\lean{Turing.FinTM.ComputesInTime.output_unique}
\label{Turing.FinTM.ComputesInTime.output_unique}
\leanok
\uses{Turing.FinTM, Turing.FinTM.ComputesInTime, Turing.FinTM.ComputesInTime.mono, Turing.MultiTapeTM.ComputesInTimeAndSpace}
\difficulty{3}
Let $M$ be a finite Turing machine over an alphabet $\Sigma$ and let $x$ be a
word over $\Sigma$.  If $M$, started on $x$, halts within $t$ steps with $w$
on its output tape and also halts within $t'$ steps with $w'$ on its output
tape, then $w = w'$.  In other words, a machine has at most one completed
output on a given input.
\end{theorem}

\begin{theorem}[Time-bounded computation implies computation]
\lean{Turing.FinTM.ComputesFunInTime.computes}
\label{Turing.FinTM.ComputesFunInTime.computes}
\leanok
\uses{Turing.FinTM, Turing.FinTM.Computes, Turing.FinTM.ComputesFunInTime}
\difficulty{1}
Let $M$ be a finite Turing machine over an alphabet $\Sigma$.  If $M$ computes
the string function $f \colon \Sigma^{*} \to \Sigma^{*}$ in time
$T \colon \mathbb{N} \to \mathbb{N}$ (halting within $T(\abs{x})$ steps on
every input $x$ with $f(x)$ on its output tape), then $M$ computes $f$, i.e.\
on every input it eventually halts with the value of $f$ on its output tape.
\end{theorem}

\begin{theorem}[Space-free characterization of timed computation]
\lean{Turing.FinTM.computesInTime_iff}
\label{Turing.FinTM.computesInTime_iff}
\leanok
\uses{Turing.FinTM, Turing.FinTM.ComputesInTime, Turing.MultiTapeTM.initCfg, Turing.MultiTapeTM.runFrom}
\difficulty{2}
Let $M$ be a finite Turing machine over an alphabet $\Sigma$, let $x$ and $w$
be words over $\Sigma$, and let $t \in \mathbb{N}$.  Then $M$ computes the
output $w$ from the input $x$ within $t$ steps if and only if the
configuration reached after running $M$ for $t$ steps from its initial
configuration on $x$ is halted and carries exactly the word $w$ on its output
tape.
\end{theorem}

\begin{theorem}[Completed outputs of a total machine]
\lean{Turing.FinTM.Computes.exists_computesInTime_iff}
\label{Turing.FinTM.Computes.exists_computesInTime_iff}
\leanok
\uses{Turing.FinTM, Turing.FinTM.Computes, Turing.FinTM.ComputesInTime, Turing.FinTM.ComputesInTime.output_unique}
\difficulty{3}
Let $M$ be a finite Turing machine over an alphabet $\Sigma$ that computes the
string function $g \colon \Sigma^{*} \to \Sigma^{*}$, and let $x$ and $w$ be
words over $\Sigma$.  Then there exists a number of steps $t$ such that $M$,
started on $x$, halts within $t$ steps with $w$ on its output tape if and only
if $w = g(x)$.
\end{theorem}
